%% Evaluation design and E1--E3. Section labels follow the historical source.
\section{Evaluation}
\label{sec:eval}
We evaluate what is stored and remains accessible after a change (E1--E2,
E4), what the retrieval path selects and returns (E3, E5, E8), and what
answer diagnostics and deployment assertions establish (E6--E7, E9).
Tables~S9--S10 and S17 give the protocols and denominators. The separate
experiment source map links the descriptive campaign names to their
archives. We do not pool engineering assertions with benchmark results.

\subsection{Experimental design and measures}
\label{sec:eval-setup}

\textbf{Corpora and questions.} The routing and evidence-pack benchmark
uses the same 72 in-domain questions across ten domains in three campaigns:
the Frozen aggregate (138 documents stored as 879 parts), the Follow-up
(906 parts), and the Main comparison (seven arms on a restarted stack).
The Main comparison is the basis for the primary contrasts; the other
campaigns show how the same labels vary across corpus and index states. Each complete arm also has
four out-of-domain probes (R1--R4); two cross-domain placeholders were not
scored. Developers authored the questions and target annotations. A separate
paper corpus has 100 target papers and ten question-specific source IDs.
The historical industrial-corpus ranking diagnostic has 16 questions on
42 documents, including twenty near-duplicate telemetry reports; Q16 shares
a token with its gold filename. Tables~S9--S10 and S17 document these and
the separate multilingual listing checks.

\textbf{Measures.} The route metric checks which collection the script
selected. Containment checks for the annotated literal string anywhere in
the returned pack, even in a non-gold document. The historical ranking
analysis uses document Hit@5; the answer diagnostics check literal source
IDs, annotated terms (at least $\lceil 3n/5\rceil$ of $n$) and marker syntax.
These measures do not substitute for judged relevance or claim support
\citep{craswell-msmarco-2021,jarvelin-cumulated-gain-2002,es-ragas-2024,saad-falcon-ares-2024}.
Wilson intervals are marginal summaries; questions are paired by ID for
two-sided exact McNemar tests. Tests are exploratory and unadjusted.
Each answer condition was run once; retries belong to the same condition,
not independent trials. The historical grids used Claude Code through the
platform chat API; the four-question engine contrast in Table~S16 is a
separate instructed case.

\subsection{E1: Acquisition and parsing}
\label{sec:e1}
The storage audit found at least one stored part for each
of the 100 target papers. The five libraries contained 328 parts, including
six outside the target set. All 328 passed the recorded description check.
Neither check measures retrieval coverage. A filename-stem diagnostic found
all target identifiers, but most searches were unscoped and chunk copies
could count. A separate ten-question vector diagnostic ranked each gold
document first within its preassigned library; it did not test automatic
routing or passage support. Table~S9 preserves the audit details.

In a separate batch fixture, the warmed parser processed three
PDFs in 38.7~s. Those three parsed documents and two directly ingested
Markdown documents were all retrieved at rank one by document-specific
probes. This checks both acquisition paths on five inputs, not parser
accuracy or general retrieval recall.

\subsection{E2: Lifecycle consistency}
\label{sec:e2}
The single-document rename fixture compared three observations:
the stored metadata and vector entries showed the new name,
but the repeated query initially showed the old one. A cache-bypassing query
returned the new name, and a probe after forced reindexing confirmed it.
This distinguishes cached output from index state in one fixture; it does not
establish concurrent-write or crash-recovery safety.
An older rename test had found duplicate chunks before forced reindexing;
the two checks describe different system states. Table~S9 and the historical
inventory retain the separate reconciliation, sub-library, tagging and
graph checks.

\subsection{E3: Collection selection versus evidence retrieval}
\label{sec:routing-audit}
In the Main comparison, the scripted selector chose the labelled collection for
all 72 in-domain questions with full route cards (A0), compared with 49
with nominally name-only cards (A3). The corresponding evidence packs
contained the target string for 36 and 26 questions, respectively
(Figure~\ref{fig:routing}). The runner logged a shelf request but chose documents by vector rank, not
from shelf descriptions. A3 could fall back to descriptions without logging
that path and could search outside the benchmark libraries. Some route cards
also share title terms with the questions. A3 returned no candidate on
35 of 72 questions, including 23 of the 49 it routed correctly. A0 found
candidates on 21 of those 23. The comparison does not isolate descriptions
or shelf navigation. Earlier runs on the same questions appear separately
in Figure~\ref{fig:routing}.

\begin{figure}[!t]
\centering
\includegraphics[width=\columnwidth]{fig_latest_route}
\caption{(a) Labelled-collection selection and (b) target-string containment
for A0 and A3. Shading marks the Main comparison. All three campaigns use
the same 72 in-domain questions; they are not independent test sets.}
\label{fig:routing}
\end{figure}
